\chapter{Vectoren, rechten en vlakken in de ruimte}\label{ch:g12:space}

De meetkunde van de drie dimensies wordt berekenbaar zodra \omterm{def:g11:vect:vector}{vectoren} en
\omterm{def:g10:coordgeom:system}{coördinaten} beschikbaar zijn: rechten krijgen een parametervoorstelling,
vlakken krijgen een cartesische \omterm{def:g10:algebra:equation}{vergelijking}, en het \omterm{def:g12:space:dot}{scalair product} meet
hoeken en afstanden. Dit hoofdstuk bouwt die gereedschapskist op en lost er de
klassieke problemen over \omterm{def:g10:numbers:interunion}{doorsneden} en afstanden mee op.

\section{Vectoren in de ruimte}

\omterm{def:g11:vect:vector}{Vectoren} in de ruimte gehoorzamen aan dezelfde regels als in het vlak: ze
tellen op, ze schalen, en $\vect{AB} = \vect{CD}$ precies wanneer $ABDC$ een
parallellogram is. In een \omterm{def:g10:coordgeom:system}{assenstelsel} $(O; \vec\imath, \vec\jmath, \vec k)$
heeft elke \omterm{def:g11:vect:vector}{vector} \omterm{def:g10:coordgeom:system}{coördinaten} $(x, y, z)$.

\begin{definition}[Collineair, coplanair]\label{def:g12:space:collinear}
Twee \omterm{def:g11:vect:vector}{vectoren} heten \emph{\omterm{def:g11:vect:collinear}{collineair}} als de ene een veelvoud van de andere
is. Drie \omterm{def:g11:vect:vector}{vectoren} $\vec u, \vec v, \vec w$ heten
\emph{coplanair}\index{coplanair} als een van hen als een combinatie van de
twee andere geschreven kan worden, bijvoorbeeld
$\vec w = a\vec u + b\vec v$.
\end{definition}

\begin{proposition}\label{prop:g12:space:basis}
Zijn $\vec u, \vec v, \vec w$ niet \omterm{def:g12:space:collinear}{coplanair}, dan valt elke \omterm{def:g11:vect:vector}{vector} van de
ruimte op precies één manier uiteen als $x\vec u + y\vec v + z\vec w$.
\end{proposition}

\begin{proof}[Idee van het bewijs]
Trek door het uiteinde van de te \omterm{def:g10:algebra:expand}{ontbinden} \omterm{def:g11:vect:vector}{vector} de rechte evenwijdig met
$\vec w$; die snijdt het vlak van $\vec u$ en $\vec v$ in één enkel punt, dat
de \omterm{def:g11:vect:vector}{vector} splitst in een component in dat vlak (op één manier
$x\vec u + y\vec v$, vlakke meetkunde) en een component langs $\vec w$.
Uniciteit: een \omterm{def:g11:seq:arithmetic}{verschil} van twee ontbindingen zou een van de \omterm{def:g11:vect:vector}{vectoren} in
\omterm{def:g11:func:function}{functie} van de twee andere uitdrukken, in tegenspraak met de
niet-coplanariteit.
\end{proof}

\section{Rechten en vlakken}

\begin{definition}[Parametervoorstelling van een rechte]\label{def:g12:space:line}
De rechte door $A(x_A, y_A, z_A)$ met \omterm{def:g11:vect:direction}{richtingsvector}
$\vec u (a, b, c) \neq \vec 0$ is de verzameling van de punten
\[
M(x_A + ta,\; y_A + tb,\; z_A + tc), \qquad t \in \R .
\]
\end{definition}

\begin{definition}[Vlak]\label{def:g12:space:plane}
Het vlak door $A$ opgespannen door twee niet-collineaire \omterm{def:g11:vect:vector}{vectoren}
$\vec u, \vec v$ is de verzameling van de punten $M$ met
$\vect{AM} = s\vec u + t\vec v$, $(s, t) \in \R^2$.
\end{definition}

\begin{omfigure}
\begin{tikzpicture}[scale=1.05]
  \draw (-0.5,-0.25) -- (3.4,1.7);
  \draw[-Stealth, omDef, thick] (1,0.5) -- (1.9,0.95)
    node[midway, above left, font=\small] {$\vec u$};
  \fill (0.1,0.05) circle (1.2pt) node[below right, font=\small] {$t=-1$};
  \fill (1,0.5) circle (1.2pt) node[above left, font=\small] {$A$};
  \fill (1.9,0.95) circle (1.2pt) node[below right, font=\small] {$t=1$};
  \fill (2.8,1.4) circle (1.2pt) node[below right, font=\small] {$t=2$};
\end{tikzpicture}
\qquad\quad
\begin{tikzpicture}[scale=1.05]
  \draw (0,0) -- (4,0.5) -- (5.2,1.9) -- (1.2,1.4) -- cycle;
  \node[font=\small] at (4.75,1.6) {$\mathcal P$};
  \coordinate (A) at (1.3,0.55);
  \coordinate (M) at (3.77,1.38);
  \draw[-Stealth, omDef, thick] (A) -- ++(1.31,0.16)
    node[midway, below, font=\small] {$\vec u$};
  \draw[-Stealth, omDef, thick] (A) -- ++(0.42,0.49)
    node[midway, left, font=\small] {$\vec v$};
  \draw[black, dashed, thin] (3.27,0.79) -- (M) -- (1.80,1.14);
  \fill (A) circle (1.2pt) node[below left, font=\small] {$A$};
  \fill (M) circle (1.2pt) node[above, font=\small] {$M$};
\end{tikzpicture}

{\small Links: de rechte door $A$ met richting $\vec u$, geijkt door de
parameter $t$. Rechts: het vlak door $A$ opgespannen door $\vec u$ en
$\vec v$; elk punt $M$ van het vlak wordt bereikt als
$\vect{AM} = s\vec u + t\vec v$.}
\end{omfigure}

\begin{proposition}[Onderlinge liggingen]\label{prop:g12:space:positions}
Twee verschillende vlakken zijn ofwel evenwijdig, ofwel snijden ze elkaar
volgens een rechte. Een rechte en een vlak zijn ofwel evenwijdig (eventueel
met de rechte in het vlak), ofwel snijden ze elkaar in één enkel punt. Twee
rechten in de ruimte kunnen snijdend, strikt evenwijdig, samenvallend of
\emph{kruisend} (niet \omterm{def:g12:space:collinear}{coplanair}) zijn.
\end{proposition}

\begin{proof}[Schets]
Het gaat om uitspraken over incidentie; elke gevalsonderscheiding herleidt
zich in \omterm{def:g10:coordgeom:system}{coördinaten} tot het oplossen van een \omterm{def:g10:lines:system}{lineair stelsel} en het tellen van
de oplossingen --- zie \cref{met:g12:space:intersections}.
\end{proof}

\section{Scalair product}

\begin{definition}[Scalair product]\label{def:g12:space:dot}
Het \emph{scalair product}\index{scalair product} van $\vec u$ en $\vec v$ is
\[
\vec u \cdot \vec v = \tfrac12\left(\norm{\vec u + \vec v}^2
- \norm{\vec u}^2 - \norm{\vec v}^2\right),
\]
waarbij $\norm{\vec u}$ de lengte van $\vec u$ is. Zijn beide \omterm{def:g11:vect:vector}{vectoren}
verschillend van $\vec 0$, dan is
$\vec u \cdot \vec v = \norm{\vec u}\,\norm{\vec v}\cos\theta$ met $\theta$ de
hoek tussen beide; en in een \emph{\omterm{def:g10:coordgeom:system}{orthonormaal}} \omterm{def:g10:coordgeom:system}{assenstelsel} geldt
\[
\vec u \cdot \vec v = xx' + yy' + zz' .
\]
Twee \omterm{def:g11:vect:vector}{vectoren} heten \emph{\omterm{def:g11:scal:orthogonal}{orthogonaal}} als hun scalair product $0$ is.
\end{definition}

\begin{proposition}[Eigenschappen]\label{prop:g12:space:dotprops}
Het \omterm{def:g12:space:dot}{scalair product} is symmetrisch ($\vec u\cdot\vec v = \vec v\cdot\vec u$)
en bilineair ($(a\vec u + b\vec v)\cdot \vec w = a\,\vec u\cdot\vec w +
b\,\vec v\cdot \vec w$), en $\vec u \cdot \vec u = \norm{\vec u}^2 \geq 0$.
\end{proposition}

\begin{proof}
In een \omterm{def:g10:coordgeom:system}{orthonormaal assenstelsel} zijn alle drie de eigenschappen onmiddellijk
af te lezen op de formule $xx' + yy' + zz'$, die zelf volgt uit de
definiërende formule en uit de berekening van de lengtes met Pythagoras:
$\norm{\vec u}^2 = x^2 + y^2 + z^2$.
\end{proof}

\begin{definition}[Normaalvector]\label{def:g12:space:normal}
Een \emph{normaalvector}\index{normaalvector} van een vlak $\mathcal P$ is een
\omterm{def:g11:vect:vector}{vector} verschillend van $\vec 0$ die \omterm{def:g11:scal:orthogonal}{orthogonaal} is met elke \omterm{def:g11:vect:vector}{vector} die in
$\mathcal P$ ligt (het volstaat dat hij \omterm{def:g11:scal:orthogonal}{orthogonaal} is met twee
niet-collineaire richtingen van $\mathcal P$).
\end{definition}

\begin{omfigure}
\begin{tikzpicture}[scale=1.1]
  \draw (0,0) -- (4,0.5) -- (5.2,1.9) -- (1.2,1.4) -- cycle;
  \node[font=\small] at (4.75,1.6) {$\mathcal P$};
  \coordinate (H) at (2.6,0.95);
  \coordinate (M0) at (2.86,3.03);
  \draw[-Stealth, omThm, thick] (H) -- (2.73,1.99)
    node[midway, right=1pt, font=\small] {$\vec n$};
  \draw[black, dashed, thin] (2.73,1.99) -- (M0);
  \draw[thin] (2.79,0.974) -- (2.815,1.174) -- (2.625,1.15);
  \draw[black, dashed, thin] (H) -- (1.1,0.55);
  \fill (H) circle (1.2pt) node[below right, font=\small] {$H$};
  \fill (M0) circle (1.2pt) node[above, font=\small] {$M_0$};
  \fill (1.1,0.55) circle (1.2pt) node[below left, font=\small] {$A$};
\end{tikzpicture}

{\small Een \omterm{def:g12:space:normal}{normaalvector} $\vec n$ van het vlak $\mathcal P$ is \omterm{def:g11:scal:orthogonal}{orthogonaal}
met elke richting van $\mathcal P$. Het punt $H$, de voet van de loodlijn uit
$M_0$, verwezenlijkt de afstand van $M_0$ tot het vlak.}
\end{omfigure}

\begin{theorem}[Cartesische vergelijking van een vlak]\label{thm:g12:space:planeeq}
In een \omterm{def:g10:coordgeom:system}{orthonormaal assenstelsel} heeft het vlak door $A$ met \omterm{def:g12:space:normal}{normaalvector}
$\vec n(a, b, c)$ als \omterm{def:g10:algebra:equation}{vergelijking}
\[
a x + b y + c z + d = 0,
\]
met $d = -(ax_A + by_A + cz_A)$. Omgekeerd bepaalt elke zulke \omterm{def:g10:algebra:equation}{vergelijking} met
$(a,b,c) \neq (0,0,0)$ een vlak met \omterm{def:g12:space:normal}{normaalvector} $(a, b, c)$.
\end{theorem}

\begin{proof}
$M \in \mathcal P \iff \vect{AM} \perp \vec n \iff
a(x - x_A) + b(y - y_A) + c(z - z_A) = 0$, wat uitgewerkt de \omterm{def:g10:algebra:equation}{vergelijking}
oplevert. Omgekeerd: kies een willekeurig punt $A$ dat aan de \omterm{def:g10:algebra:equation}{vergelijking}
voldoet; dezelfde berekening in omgekeerde richting toont dat de
oplossingsverzameling $\{M : \vect{AM}\cdot\vec n = 0\}$ is, het vlak door $A$
loodrecht op $\vec n$.
\end{proof}

\begin{method}[Doorsneden in de praktijk]\label{met:g12:space:intersections}
\begin{itemize}
  \item \emph{Rechte $\cap$ vlak}: vul de \omterm{def:g10:coordgeom:system}{coördinaten} uit de
        parametervoorstelling van de rechte in de \omterm{def:g10:algebra:equation}{vergelijking} van het vlak in
        en los op naar $t$ (één oplossing: één punt; $0 = $ iets van nul
        verschillend: evenwijdig; $0 = 0$: de rechte ligt in het vlak).
  \item \emph{Vlak $\cap$ vlak}: los het stelsel van de twee \omterm{def:g10:algebra:equation}{vergelijkingen} op
        en parametriseer met één vrije coördinaat; de oplossing is een rechte
        (of de vlakken zijn evenwijdig wanneer de \omterm{def:g12:space:normal}{normaalvectoren} \omterm{def:g11:vect:collinear}{collineair}
        zijn).
  \item \emph{Toetsen op orthogonaliteit}: een rechte staat loodrecht op een
        vlak als en slechts als haar richting \omterm{def:g11:vect:collinear}{collineair} is met de normaal;
        twee vlakken staan loodrecht op elkaar als en slechts als hun normalen
        \omterm{def:g11:scal:orthogonal}{orthogonaal} zijn.
\end{itemize}
\end{method}

\begin{proposition}[Afstand van een punt tot een vlak]\label{prop:g12:space:distance}
\index{afstand!punt tot vlak}
In een \omterm{def:g10:coordgeom:system}{orthonormaal assenstelsel} is de afstand van $M_0(x_0, y_0, z_0)$ tot
het vlak $\mathcal P : ax + by + cz + d = 0$ gelijk aan
\[
\operatorname{dist}(M_0, \mathcal P)
= \frac{\abs{ax_0 + by_0 + cz_0 + d}}{\sqrt{a^2 + b^2 + c^2}}.
\]
\end{proposition}

\begin{proof}
Zij $H$ de \omterm{def:g11:scal:orthogonal}{orthogonale} projectie van $M_0$ op $\mathcal P$: de voet van de
loodlijn, zodat $\vect{HM_0}$ \omterm{def:g11:vect:collinear}{collineair} is met de eenheidsnormaal
$\frac{\vec n}{\norm{\vec n}}$ en de afstand
$\abs{\vect{HM_0}\cdot \frac{\vec n}{\norm{\vec n}}}$ bedraagt. Voor elke
$H \in \mathcal P$ geldt
$\vect{HM_0}\cdot\vec n = ax_0 + by_0 + cz_0 - (ax_H + by_H + cz_H)
= ax_0 + by_0 + cz_0 + d$ (met de \omterm{def:g10:algebra:equation}{vergelijking} in $H$), waaruit de formule
volgt na deling door $\norm{\vec n} = \sqrt{a^2+b^2+c^2}$.
\end{proof}

\begin{example}\label{ex:g12:space:distance}
Afstand van de oorsprong tot het vlak $x + 2y + 2z - 6 = 0$:
$\frac{\abs{-6}}{\sqrt{1 + 4 + 4}} = \frac63 = 2$.
\end{example}

\section{Oefeningen}

In alle oefeningen is het \omterm{def:g10:coordgeom:system}{assenstelsel} \omterm{def:g10:coordgeom:system}{orthonormaal}.

\begin{exercise}[$\star$]\label{exo:g12:space:1}
Zij $A(1, 0, 2)$, $B(3, 1, 1)$ en $C(2, -1, 3)$. Liggen de punten $A$, $B$ en
$C$ op één rechte? Geef een parametervoorstelling van de rechte $(AB)$.
\end{exercise}

\begin{exercise}[$\star$]\label{exo:g12:space:2}
Bepaal een cartesische \omterm{def:g10:algebra:equation}{vergelijking} van het vlak door $A(1, 1, 0)$ met
\omterm{def:g12:space:normal}{normaalvector} $\vec n(2, -1, 3)$. Behoort het punt $B(0, 2, 1)$ ertoe?
\end{exercise}

\begin{exercise}[$\star$]\label{exo:g12:space:3}
Bereken de hoek in $A$ (op één graad nauwkeurig) van de driehoek met
hoekpunten $A(0,0,0)$, $B(1, 1, 0)$ en $C(1, 0, 1)$.
\end{exercise}

\begin{exercise}[$\star\star$]\label{exo:g12:space:4}
Beschouw de rechte $\mathcal D$ door $A(1, 2, 0)$ met richting
$\vec u(1, -1, 2)$, en het vlak $\mathcal P : 2x + y - z + 1 = 0$. Bepaal
$\mathcal D \cap \mathcal P$.
\end{exercise}

\begin{exercise}[$\star\star$]\label{exo:g12:space:5}
Zij $\mathcal P : x - y + 2z = 1$ en $\mathcal Q : 2x + y + z = 4$. Toon aan
dat $\mathcal P$ en $\mathcal Q$ elkaar volgens een rechte snijden en geef
daarvan een parametervoorstelling.
\end{exercise}

\begin{exercise}[$\star\star$]\label{exo:g12:space:6}
Toon aan dat de rechten
\[
\mathcal D_1 : (x, y, z) = (1 + t,\; 2t,\; 3 - t), \qquad
\mathcal D_2 : (x, y, z) = (2 + s,\; 1 + s,\; 1 + 2s)
\]
kruisend zijn (niet \omterm{def:g12:space:collinear}{coplanair}).
\end{exercise}

\begin{exercise}[$\star\star$]\label{exo:g12:space:7}
Zij $\mathcal P : 2x - y + 2z - 3 = 0$.
\begin{enumerate}
  \item Bereken de afstand van $M_0(3, 1, 1)$ tot $\mathcal P$.
  \item Bepaal de \omterm{def:g10:coordgeom:system}{coördinaten} van de \omterm{def:g11:scal:orthogonal}{orthogonale} projectie $H$ van $M_0$ op
        $\mathcal P$, en controleer de afstand $M_0H$.
\end{enumerate}
\end{exercise}

\begin{exercise}[$\star\star\star$]\label{exo:g12:space:8}
De kubus $ABCDEFGH$ heeft ribbe $1$; kies de \omterm{def:g10:coordgeom:system}{coördinaten} zo dat
$A(0,0,0)$, $B(1,0,0)$, $D(0,1,0)$ en $E(0,0,1)$ (met
$C = B + D$, $F = B + E$, $G = B + D + E$ en $H = D + E$ als \omterm{def:g11:vect:vector}{vectoren} vanuit
$A$).
\begin{enumerate}
  \item Toon aan dat de diagonaal $(AG)$ \omterm{def:g11:scal:orthogonal}{orthogonaal} is met het vlak $(BDE)$.
  \item Bereken de afstand van $A$ tot het vlak $(BDE)$, en het punt waar
        $(AG)$ er doorheen prikt.
\end{enumerate}
\end{exercise}

\begin{exercise}[$\star\star\star$]\label{exo:g12:space:9}
\emph{(Volume van een viervlak.)} Zij $A(1,1,1)$, $B(2,1,0)$, $C(0,2,1)$ en
$D(2,2,2)$.
\begin{enumerate}
  \item Ga na dat $\vect{AB}$, $\vect{AC}$ en $\vect{AD}$ niet \omterm{def:g12:space:collinear}{coplanair} zijn
        (de punten vormen een echt viervlak).
  \item Zoek een \omterm{def:g11:vect:vector}{vector} $\vec n$ \omterm{def:g11:scal:orthogonal}{orthogonaal} met zowel $\vect{BC}$ als
        $\vect{BD}$, en leid een cartesische \omterm{def:g10:algebra:equation}{vergelijking} van het vlak $(BCD)$
        af.
  \item Bereken de afstand van $A$ tot het vlak $(BCD)$ en de oppervlakte van
        de driehoek $BCD$, met de formule
        $\text{Opp} = \frac12\sqrt{\norm{\vec u}^2\norm{\vec v}^2 -
        (\vec u \cdot \vec v)^2}$ voor de driehoek opgespannen door
        $\vec u$ en $\vec v$.
  \item Leid het volume van het viervlak af
        ($V = \frac13 \times \text{grondvlak} \times \text{hoogte}$).
\end{enumerate}
\end{exercise}

\section{Opgave: de kubus doorsnijden}

\begin{problem}[{Weekendopgave --- een vlak snijdt een kubus tot een
volmaakte zeshoek, een regelmatig viervlak verstopt zich in de hoeken, en de
hoek van de diagonaal vormt elke diamant}]
\label{pb:g12:space:1}
Snijd een kubus door en je verwacht vierkanten en rechthoeken --- en toch
levert één beroemde snede een \emph{regelmatige zeshoek} op, en laten vier
slimme hoeksnedes een \emph{regelmatig viervlak} achter. Beide verrassingen,
en de hoek van $109.5^\circ$ die de koolstofatomen vereren, bezwijken voor de
gereedschapskist van dit hoofdstuk: parametervoorstellingen van rechten,
\omterm{def:g10:algebra:equation}{vergelijkingen} van vlakken en het \omterm{def:g12:space:dot}{scalair product}
(\cref{thm:g12:space:planeeq}, \cref{prop:g12:space:distance}). Gebruik overal
de kubus van \cref{exo:g12:space:8}: ribbe $1$, $A(0,0,0)$, $B(1,0,0)$,
$D(0,1,0)$, $E(0,0,1)$, $C(1,1,0)$, $F(1,0,1)$, $H(0,1,1)$, $G(1,1,1)$.

\medskip
\noindent\textbf{Deel I --- Vlotheid.}
\begin{enumerate}
  \item Schrijf de \omterm{def:g10:algebra:equation}{vergelijking} op van het vlak door $B(1,0,0)$ met
        \omterm{def:g12:space:normal}{normaalvector} $(1,1,1)$, een parametervoorstelling van de rechte door
        $A$ met richting $(1,1,1)$, en hun snijpunt.
  \item Bereken de afstand van $A$ tot dat vlak
        (\cref{prop:g12:space:distance}).
  \item Zijn de \omterm{def:g11:vect:vector}{vectoren} $(1,-1,0)$ en $(1,1,-2)$ \omterm{def:g11:scal:orthogonal}{orthogonaal}?
  \item Hoe liggen de vlakken $x + y + z = 1$ en $2x + 2y + 2z = 5$ ten
        opzichte van elkaar? En de rechte door $A$ met richting $(1,1,0)$ ten
        opzichte van het eerste vlak?
  \item Som de \omterm{prop:g10:coordgeom:midpoint}{middens} op van de zes ribben $[BC], [CD], [DH], [HE], [EF],
        [FB]$ van de kubus, en ga na dat alle zes aan de \omterm{def:g10:algebra:equation}{vergelijking}
        $x + y + z = \frac32$ voldoen.
\end{enumerate}

\medskip
\noindent\textbf{Deel II --- De zeshoekige \omterm{def:g10:numbers:interunion}{doorsnede}.}
Zij $P$ het vlak $x + y + z = \frac32$.
\begin{enumerate}[resume]
  \item Toon aan dat $P$ loodrecht staat op de lichaamsdiagonaal $(AG)$ en
        door het middelpunt van de kubus gaat.
  \item Bereken de afstanden tussen de opeenvolgende \omterm{prop:g10:coordgeom:midpoint}{middens} uit vraag 5 langs
        de \omterm{def:g10:numbers:interunion}{doorsnede}: wat vind je?
  \item Bereken de hoek van de veelhoek van de \omterm{def:g10:numbers:interunion}{doorsnede} in één hoekpunt (twee
        opeenvolgende zijdevectoren en een \omterm{def:g12:space:dot}{scalair product}). Besluit: de
        \omterm{def:g10:numbers:interunion}{doorsnede} is een \emph{regelmatige zeshoek}.
  \item Bereken de omtrek en de oppervlakte van de zeshoek, en vergelijk die
        oppervlakte met die van een zijvlak van de kubus. (Voor de
        volledigheid: de grootste vlakke \omterm{def:g10:numbers:interunion}{doorsnede} van de eenheidskubus is de
        diagonale rechthoek $1 \times \sqrt2$, met oppervlakte $\sqrt2$ --- de
        regelmatige zeshoek pakt het zilver.)
  \item Waar prikt de diagonaal $(AG)$ door $P$? Ga na dat het prikpunt het
        \omterm{prop:g10:coordgeom:midpoint}{midden} van $[AG]$ is.
  \item Laat de snede schuiven: beschrijf de \omterm{def:g10:numbers:interunion}{doorsneden} $x + y + z = c$
        naarmate $c$ van $0$ tot $3$ groeit --- bereken de \omterm{def:g10:numbers:interunion}{doorsnede} voor
        $c = \frac12$ (welke veelhoek, welke afmeting?), en vertel de
        gedaanteverwisseling die bij $c = \frac32$ door jouw zeshoek gaat.
  \item Waarom kan geen enkel vlak een kubus in een veelhoek met \emph{zeven}
        zijden snijden? (Waar moet elke zijde van de \omterm{def:g10:numbers:interunion}{doorsnede} liggen?)
\end{enumerate}

\medskip
\noindent\textbf{Deel III --- Het verborgen viervlak.}
\begin{enumerate}[resume]
  \item Bereken het volume van het hoekviervlak $ABDE$ (grondvlak, hoogte en
        de formule met een derde die het onderbouwvolume aannam ---
        \cref{exo:g12:space:9} berekent zulke volumes in het algemeen).
  \item De vier hoekpunten $B$, $D$, $E$ en $G$: bereken alle zes de
        paarsgewijze afstanden en besluit dat $BDEG$ een \emph{regelmatig}
        viervlak is. Bepaal daarna zijn volume door de hoekviervlakken van de
        kubus af te trekken.
  \item Bereken de afstand van het middelpunt van de kubus tot het vlak
        $(BDE)$.
  \item De diamanthoek: bereken de hoek tussen de diagonalen $\vect{AG}$ en
        $\vect{BH}$ van de kubus. Haar supplement, ongeveer $109.5^\circ$, is
        de hoek tussen de bindingen in diamant en in methaan --- waarom kiezen
        vier elektronenparen rond een koolstofatoom de hoekpunten van het
        viervlak uit vraag 14?
\end{enumerate}

\medskip
\noindent\textbf{Deel IV --- Echt driedimensionaal.}
\begin{enumerate}[resume]
  \item Zoek waar de rechte door $B(1,0,0)$ met richting $(0,1,1)$ door het
        vlak $P$ van de zeshoek prikt.
  \item Toon aan dat de rechten $(AB)$ en $(EG)$ elkaar niet snijden en niet
        evenwijdig zijn: \emph{kruisende} rechten, het verschijnsel dat in een
        vlak onmogelijk is.
  \item Bereken de \omterm{def:g11:scal:orthogonal}{orthogonale} projectie van $B$ op $P$ (loop van $B$ uit
        langs de normaal tot aan het vlak), en ga na dat jouw punt aan de
        \omterm{def:g10:algebra:equation}{vergelijking} van $P$ voldoet.
  \item Slotstuk --- de uitrusting van de landmeter in de ruimte:
        parametervoorstellingen om door te prikken, \omterm{def:g12:space:normal}{normaalvectoren} voor
        hoeken en afstanden, \omterm{def:g10:algebra:equation}{vergelijkingen} van vlakken voor \omterm{def:g10:numbers:interunion}{doorsneden}; en de
        kubus als speelterrein waar ze een zeshoek, een regelmatig viervlak,
        de diamanthoek en een paar kruisende rechten voortbrachten. Telkens
        één zin.
\end{enumerate}
\end{problem}
